% Standalone chunk A.7: Hard-projection regret lower bound
% (POST-FIX, Round 2).
% Covers: Thm spsc_lb statement (Pi_HP class, large-d regime),
%         Le Cam proof (two-point, adversarial-v, exploit-KL suppression),
%         Rem lb_scope_conf, Rem lb_class_conf.
% Source: conf.tex lines ~516-541, 1947-2077 (post-fix).
% Shared preamble for standalone chunks. Do not compile directly; \input{} it.
\documentclass{article}
\usepackage[utf8]{inputenc}
\usepackage[T1]{fontenc}
\usepackage{lmodern}
\usepackage[margin=1in]{geometry}
\usepackage{amsmath,amssymb,amsfonts,amsthm,bm,mathtools}
\usepackage{enumitem}
\usepackage{microtype}
\usepackage{xcolor}
\usepackage{natbib}
\usepackage{booktabs}
\usepackage{hyperref}

\newcommand{\R}{\mathbb{R}}
\newcommand{\N}{\mathbb{N}}
\newcommand{\E}{\mathbb{E}}
\newcommand{\Prob}{\mathbb{P}}
\newcommand{\op}{\mathrm{op}}
\DeclareMathOperator{\sign}{sign}
\newcommand{\cA}{\mathcal{A}}
\newcommand{\cF}{\mathcal{F}}
\newcommand{\cH}{\mathcal{H}}
\newcommand{\cT}{\mathcal{T}}
\newcommand{\cI}{\mathcal{I}}
\newcommand{\cW}{\mathcal{W}}
\newcommand{\cTprobe}{\mathcal{T}_{\mathrm{probe}}}
\newcommand{\cE}{\mathcal{E}}
\newcommand{\norm}[1]{\left\lVert #1 \right\rVert}
\newcommand{\tilO}{\widetilde{\mathcal{O}}}
\newcommand{\DynReg}{\mathrm{DynReg}}
\newcommand{\rank}{\mathrm{rank}}
\newcommand{\range}{\mathrm{range}}
\DeclareMathOperator{\TopEig}{TopEig}
\DeclareMathOperator*{\argmax}{arg\,max}
\DeclareMathOperator{\tr}{tr}

\newtheorem{assumption}{Assumption}
\newtheorem{theorem}{Theorem}
\newtheorem{corollary}{Corollary}
\newtheorem{proposition}{Proposition}
\newtheorem{lemma}{Lemma}
\newtheorem{remark}{Remark}


\title{Chunk A.7: Hard-projection regret lower bound\\(post-fix, Round 2 verification)}
\author{(excerpt for standalone review)}
\date{}

\begin{document}
\maketitle

\textbf{Setup context.}
\begin{itemize}[leftmargin=*,itemsep=1pt,topsep=1pt]
\item $\Pi_{\mathrm{HP}}$: policies whose exploitation action factors
through a data-measurable rank-$r$ projector,
$x_t=\widehat P_t x_t$ (contains SPSC and, to our knowledge, every
prior algorithm for this problem).
\item Costed dynamic regret
$\DynReg_T^{(c)} := \sum_t(\max_x x^\top\theta_t-x_t^\top\theta_t)+c|\cTprobe|$.
\end{itemize}

\textbf{What changed since Round~1 (delta highlights).}
\begin{itemize}[leftmargin=*,itemsep=1pt,topsep=1pt]
\item \emph{Linear-in-$\alpha$ per-round regret.} Rewritten using
$x_t^\top v \le R_\cA\|\widehat P_t v\|$ (from $x_t=\widehat P_t x_t$),
avoiding the spurious $\alpha^2/\|w\|$ term that appeared when the
original sketch assumed $w\ne 0$.
\item \emph{Constants explicit in Le Cam.} $\alpha^2 :=
\sigma_\varepsilon^2/\bigl[4(m+R_\cA^2(T-m)r/d)\bigr]$ gives
$\mathrm{KL}\le 1/8$, so Pinsker $d_{\mathrm{TV}}\le\sqrt{1/16}=1/4$.
\item \emph{Scope-regime remarks.} Rem.~\ref{rem:lb_scope_conf} pins
down the large-$d$ regime $d\ge d_0\cdot r\,T^{1/3}$ (Pinsker bound
degrades to $\Omega(\sqrt T)$ in $d=O(r)$); Rem.~\ref{rem:lb_class_conf}
makes explicit where the $\Pi_{\mathrm{HP}}$ restriction enters.
\end{itemize}

\section*{Statement}

\begin{theorem}[Lower bound within $\Pi_{\mathrm{HP}}$]
\label{thm:spsc_lb}
Assume $d \ge d_0\cdot r\,T^{1/3}$ for an absolute constant $d_0$
(large-$d$ / subspace-recovery-limited regime). For all sufficiently
large $T$, there exists $c_0 > 0$ (depending only on
$R_\cA, \sigma_\varepsilon, d_0$) such that
\[
\inf_{\pi\in\Pi_{\mathrm{HP}}} \sup_\nu\ \E^\pi_\nu[\DynReg_T^{(c)}]
 \;\ge\; c_0\,c^{1/3}\,(R_\cA\sigma_\varepsilon)^{2/3}\,T^{2/3}.
\]
Combined with Thm.\ spsc\_regret, SPSC is rate-optimal within
$\Pi_{\mathrm{HP}}$ in this regime.
\end{theorem}

\section*{Proof}

Two-point Le Cam: $\nu_0$ has zero signal ($\theta^{(0)}\equiv 0$),
$\nu_1(v)$ has signal $\alpha v$ for a unit $v$ chosen by the
adversary after seeing the policy. Under $\Pi_{\mathrm{HP}}$,
adversarial $v$ suppresses exploit-KL, forcing identification via
probes.

\emph{Construction.} $B^{(0)}=[e_1,\dots,e_r]$, $\theta^{(0)}\equiv 0$
(limit $\Sigma_\eta\to 0^+$, $w_0=0$; bound is continuous in $\Sigma_\eta$).
$B^{(1)}=[e_1,\dots,e_{r-1},v]$, $\theta^{(1)}\equiv\alpha v$. Both
valid rank-$r$. $\varepsilon_t\sim\mathcal N(0,\sigma_\varepsilon^2)$,
$\cA=\{\|x\|\le R_\cA\}$, probes $u_t\sim\mathcal N(0,I_d)$.

\emph{Step 1 (linear-in-$\alpha$ per-round regret).} Under $\pi\in\Pi_{\mathrm{HP}}$,
$x_t=\widehat P_t x_t$, so $x_t^\top v = x_t^\top\widehat P_t v\le R_\cA\|\widehat P_t v\|$
(Cauchy--Schwarz). Optimal $\nu_1$-reward is $R_\cA\alpha$ (play $R_\cA v$),
so
\[
\Delta_t^{(1)}\ge R_\cA\alpha(1-\|\widehat P_t v\|),
\]
linear in $\alpha$. Under $\nu_0$, control regret is $0$ and costed
regret is the probe cost $cm$.

\emph{Step 2 (adversarial $v$ suppresses exploit-KL).} Sequential KL
chain rule:
\[
\mathrm{KL}(\mathbb P_0^{\cH}\|\mathbb P_1^{\cH})
=\frac{\alpha^2}{2\sigma_\varepsilon^2}\E_{\mathbb P_0}\!\Bigl[\sum_{t=1}^T(x_t^\top v)^2\Bigr].
\]
Probes ($x_t=u_t\sim\mathcal N(0,I_d)$) contribute $\E[(u^\top v)^2]=1$
each, total $m\alpha^2/(2\sigma_\varepsilon^2)$. Exploit rounds:
$(x_t^\top v)^2=(x_t^\top\widehat P_t v)^2\le R_\cA^2\|\widehat P_t v\|^2$.

Under $\nu_0$, history is noise-only ($y_t\sim\mathcal N(0,\sigma_\varepsilon^2)$,
independent of $v$), so law of $\widehat P_t$ under $\nu_0$ does not
depend on $v$. Adversary picks $v^\star$ \emph{after} this law is
determined:
\[
v^\star:=\arg\min_{\|v\|=1}\E_{\mathbb P_0(v)}\!\Bigl[\sum_{t\in\cT_{\mathrm{exp}}}\|\widehat P_t v\|^2\Bigr].
\]
Average over $v$ uniform on $\mathbb S^{d-1}$: $\E_v\|\widehat P_t v\|^2=\tr(\widehat P_t)/d=r/d$.
So $\min_v\le r/d$ per round, total $\le(T-m)r/d$. Substituting,
\[
\mathrm{KL}(\mathbb P_0^{\cH}\|\mathbb P_1^{\cH}(v^\star))
\le\frac{\alpha^2}{2\sigma_\varepsilon^2}\Bigl(m+\frac{R_\cA^2(T-m)r}{d}\Bigr).
\]

\emph{Step 3 (Le Cam).} Pinsker: $d_{\mathrm{TV}}\le\sqrt{\mathrm{KL}/2}$.
Set
\[
\alpha^2 := \frac{\sigma_\varepsilon^2}{4\bigl(m+R_\cA^2(T-m)r/d\bigr)}
\implies\mathrm{KL}\le 1/8,\quad d_{\mathrm{TV}}\le\sqrt{1/16}=1/4 < 1/2.
\]
Two-point Le Cam:
$\max(\E^\pi_{\nu_0}[\mathrm{Reg}],\E^\pi_{\nu_1(v^\star)}[\mathrm{Reg}])
\ge\tfrac12(1-d_{\mathrm{TV}})\cdot(\text{regret gap})$.

\emph{Step 4 (large-$d$ regime).} Under $d\ge d_0 r T^{1/3}$,
$R_\cA^2(T-m)r/d\le R_\cA^2 T^{2/3}/d_0$. For $m\ge T^{2/3}/d_0$, the
denominator in $\alpha^2$ is dominated by $m$, so
$\alpha\asymp\sigma_\varepsilon/\sqrt m$. By Step 1,
\[
\E_{\nu_1(v^\star)}[\mathrm{Reg}]
\ge R_\cA\alpha(T-m)\cdot(1-\sqrt{r/d}-d_{\mathrm{TV}})
\ge c_1 R_\cA\alpha(T-m)
\]
for an absolute constant $c_1>0$ (using $\sqrt{r/d}\le 1/\sqrt{d_0 T^{1/3}}$
small and $d_{\mathrm{TV}}\le 1/4$). Substituting,
\[
\sup_\nu\E^\pi_\nu[\mathrm{DynReg}_T^{(c)}]
\ge cm+c_1 R_\cA\sigma_\varepsilon T/\sqrt m
\]
(bounding $T-m\ge T/2$ for $m\le T/2$; for $m>T/2$ the probe cost
$cm\ge cT/2$ alone dominates). Minimising in $m$,
$m^\star=\Theta((R_\cA\sigma_\varepsilon T/c)^{2/3})$, and
\[
\inf_{\pi\in\Pi_{\mathrm{HP}}}\sup_\nu\E^\pi_\nu[\mathrm{DynReg}_T^{(c)}]
\ge c_0\,c^{1/3}(R_\cA\sigma_\varepsilon)^{2/3}T^{2/3}
\]
for an absolute $c_0>0$.\qed

\begin{remark}[Regime of the lower bound]
\label{rem:lb_scope_conf}
The $T^{2/3}$ rate is established for $d\ge d_0\cdot r\,T^{1/3}$
(i.e.\ $d/r\gtrsim T^{1/3}$), the natural regime in which subspace
recovery is the bottleneck. In the opposite regime $d=O(r)$, the
exploit-KL term $R_\cA^2(T-m)r/d$ is comparable to $T$, and the above
Pinsker bound degrades to $\Omega(\sqrt T)$; a tighter analysis
(e.g., instance-specific Grassmannian packing) is needed for
$d=O(r)$ but lies outside scope.
\end{remark}

\begin{remark}[Scope of the lower bound w.r.t.\ policy class]
\label{rem:lb_class_conf}
The hard-projection restriction enters via Step 2's inequality
$|x_t^\top v|=|x_t^\top\widehat P_t v|\le R_\cA\|\widehat P_t v\|$,
valid because $x_t=\widehat P_t x_t$. A policy outside
$\Pi_{\mathrm{HP}}$ could pick exploitation actions with non-zero
components along $v$ even when $\widehat P_t v=0$, contributing extra
exploit-KL that might beat $T^{2/3}$ by exploiting identifying
information from off-subspace exploitation. Constructing (or ruling
out) such an algorithm is the open problem our analysis isolates.
\end{remark}

\end{document}
